
Large language models are deployed in applications where factual accuracy is critical, yet they produce errors ranging from uncertain guessing to confident hallucination. Practitioners deploying uncertainty estimation face a landscape of methods validated independently under different conditions: semantic entropy (Kuhn et al., 2023) measures diversity across semantically equivalent outputs, self-consistency (Wang et al., 2022) checks agreement across samples, verbalized confidence (Kadavath et al., 2022) elicits self-reported uncertainty, and token probability methods analyze distributional properties. Without systematic comparison under controlled conditions, method selection lacks empirical foundation.

The research question is whether these methods capture distinct uncertainty dimensions or measure the same underlying signal through different computations. If methods are redundant, the field should consolidate approaches. If methods measure orthogonal properties, practitioners could combine signals for robust uncertainty quantification. A related question is whether different error types—knowledge gaps where models lack information versus confident misconceptions where models believe false information—exhibit characteristic uncertainty signatures that could guide method selection.

We address these questions through controlled experiments designed to isolate specific mechanisms. For the semantic clustering contribution, we compare semantic entropy (with clustering) to an ensemble baseline (without clustering) while holding computational budget constant at K=10 samples. This ablation design attributes performance differences to the clustering mechanism rather than to sampling cost. For method independence, we measure pairwise correlations across four methods applied to identical samples. For error-type signatures, we compare diversity and agreement distributions across two benchmarks representing different error characteristics.

Our experiments used Mistral-7B on 100-sample subsets from NaturalQuestions and TruthfulQA. While pilot-scale, the controlled design provides evidence for three claims: semantic clustering contributes 0.09 AUROC improvement beyond multiple sampling, maximum method correlation is 0.21 indicating orthogonal dimensions, and dataset-level error-type partitioning does not produce the predicted signatures.
